\section{Introduction}

Prediction models for nonlinear controlled systems support state estimation,
receding-horizon optimization, fault diagnosis, and simulation-based control.
Their accuracy depends on two coupled objects: the state supplied at the
forecast origin and the transition used to propagate that state under future
inputs. Learning-based model predictive control makes the same dependence
explicit because model error enters each optimized prediction horizon
\cite{hewing2020_learning_mpc}. Modern world models extend this idea with
recurrent, latent, or transformer transitions and have demonstrated strong
planning performance in continuous control and reinforcement learning
\cite{hafner2025_mastering_diverse_control_tasks_through,hansen2023_td_mpc2_scalable_robust_world,micheli2022_transformers_are_sample_efficient_world,zhang2023_storm_efficient_stochastic_transformer_based}.
For physical systems, transition accuracy and state-estimation reliability are
joint requirements when recent measurements are noisy or missing.

Recursive state estimation classically balances a dynamic prediction against
new measurements. The Kalman filter computes this balance from a stochastic
linear model and evolving error covariance \cite{kalman1960_filtering}, while
complementary filters and nonlinear observers exploit structural separation
between measurement channels and model dynamics
\cite{mahony2008_complementary,zhao2023_wheel_rail_observer}. Learned filters
such as KalmanNet relax parts of the assumed model, and data-assimilation
research couples machine learning with dynamical estimation and uncertainty
quantification
\cite{revach2022_kalmannet,cheng2023_machine_learning_with_data_assimilation}.
These methods motivate a direct question for multistep prediction: can an
explicit, fixed context correction improve a physics-residual transition when
observations are incomplete or corrupted?

The answer depends on the relative reliability of the physical prediction and
the measurement. Parameter drift, unmodeled friction, discretization, and
actuator dynamics bias nominal propagation. Measurement noise, missing packets,
and preprocessing errors bias or remove the observation. Last-value imputation
supplies a numerical input but does not make it current; pure nominal
propagation remains current but inherits model mismatch. Physics-guided
learning combines mechanistic and neural components through constraints,
features, hybrid solvers, and probabilistic state models
\cite{karniadakis2021_physics_informed_machine_learning,willard2022_integrating_scientific_knowledge_with_machine,yu2024_learning_dynamical_systems_from_data,shen2023_differentiable_modelling_to_unify_machine}.
The forecast-origin correction nevertheless remains a distinct design choice
from the learned transition.

We study this choice with the Complementary-Filtered Residual World Model
(\method). Within an eight-step observed context, nominal dynamics predict the
next state and a fixed gain of 0.5 combines that prediction with each available
measurement. Missing channels receive zero measurement gain. The final
corrected state initializes a probabilistic nominal-physics-plus-bounded-
residual transition trained by multistep Gaussian likelihood and rolled out
freely under known future controls. The update is a fixed-gain complementary
filter that exposes the allocation between nominal prediction and measurement
at the forecast origin.

The evaluation has two layers. First, five controlled nonlinear systems: a
damped oscillator, Duffing oscillator, controlled Van der Pol oscillator,
torque-driven pendulum, and cart--pole. They provide known equations, controlled
parameter shift, and paired observation corruption. Eleven transition variants
are compared across clean in-distribution, corrupted in-distribution, and
corrupted parameter-shift regimes. The primary analysis averages the 15
system-regime cells within each of eight independent data-training seeds and
tests a practical 5\% improvement margin against a preselected physics-guided
neural-network comparator. Fixed gains 0, 0.25, 0.5, 0.75, and 1 isolate the
context correction, while a one-step-trained variant isolates the training
horizon.

Second, an external experiment uses measured flight states from the NanoDrone
benchmark \cite{busetto2026_nanodrone}. Three flights from a held-out
trajectory family are evaluated with a reduced quadrotor prior, three
parameter-matched learned models, a nominal model, four context conditions, and
25- and 50-step horizons. Noise and missingness are varied separately and
together. This experiment tests whether the controlled-benchmark finding
persists when the physical prior is structurally incomplete and the targets are
recorded flight states.

The main contributions are:
\begin{enumerate}
 \item a mask-gated fixed-gain context correction coupled to a probabilistic
 physics-residual free rollout for nonlinear controlled dynamics;
 \item an error identity for correlated nominal-prediction and measurement
 errors, its constrained MSE-optimal scalar gain, and fixed-gain sensitivity;
 \item a paired confirmatory benchmark comprising \ConfirmatoryRuns{} trained
 models, five systems, three regimes, and eight independent seed clusters; and
 \item secondary classical-filter, residual-scale, structured-outage,
 projection, latency, and calibration audits, together with an external
 measured-flight evaluation.
\end{enumerate}

The controlled benchmark supports the aggregate primary hypothesis against the
preselected PGNN-style comparator but not the gated clean-regime noninferiority
requirement. Classical component analysis and residual-scale sensitivity show
that the gain-0.5 filter followed by nominal rollout performs better than the
learned residual transition. The flight
evaluation shows that lower long-horizon position error can coexist with higher
full-state error and undercovered intervals. These component and uncertainty
trade-offs matter when a predictor supplies an observer or predictive
controller.

The remainder of the paper reviews state-estimation, physics-guided, and world-
model approaches; defines the correction and free rollout; reports the
controlled and flight experiments; and discusses implications for learned
observers and receding-horizon prediction.
